\scanpage{111}
\mspage{345}

\deleted{and we in fact then have
\[
g\rho(g)\rho^2(g)=g^3=\rho^{-3}=-1
\]}
\inserted{$=\begin{pmatrix}0&\unsure{-1|+1}\\1&-1\end{pmatrix}$}

\inserted{Remarks} \deleted{\unsure{N.B.}} 1°) Obviously $-\rho$ is also a solution of (\unsure{87})
(just like $\rho^{-1}$) \deleted{\unsure{and $\rho$ thus} $(-\rho)\rho(g)\rho^2(g)=1$}
(just like $-\rho^{-1}$), but it falls
respectively under cases \deleted{\illegible{1}} 1°) \deleted{\unsure{and under case}}
\deleted{\unsure{A, resp.}}

2°) To \deleted{\illegible{1}}\inserted{determine} the solutions in $\mathrm{Gl}(2,\widehat{\mathbb Z})$
of $g\rho(g)\rho^2(g)=1$, it suffices to choose for
\unsure{$\forall\,\ell$} such a solution in $\mathrm{Gl}(2,\widehat{\mathbb Z})$, using
the \unsure{result|results} (86); we shall note that
the ``cases'' obtained 1°) to 3°) can vary from one
$\ell$ to another. \deleted{Same remark for
the solutions of the equation $g\rho(g)\rho^2(g)=\unsure{-1|1}$.}
We should keep in mind that cases 1°), 2°), 3°)
correspond \unsure{resp.\textsuperscript{t}} to families of
solutions depending on 2, 1 and 0 parameters.

\unsure{3°)} I do not know whether a $\xi\in\mathrm{Sl}(2,\widehat{\mathbb Z})^{\unsure{\prime|*}}$, \inserted{$\xi=\begin{pmatrix}A&B\\C&D\end{pmatrix}$}, \inserted{which satisfies}
$g\rho(g)\rho^2(g)=1$ admits a parametric
representation of the form (58)
by $\lambda_0\rho(\lambda_0)^{-1}$ or $\lambda'_0(-l_1^{-1})\rho(\lambda'_0)^{-1}$
\inserted{$=\lambda_0(-\unsure{l_1^{-1}|l_0^{-1}})\rho(\lambda_0)^{-1}$ with $\lambda_0\unsure{=\lambda'_0\tau_\infty}$}
(it suffices
to look separately in the $\mathrm{Sl}(2,\mathbb Z_\ell)'$).
I do not know even if we assume $g\equiv\mathrm{id}\ (2)$,
(which is a condition on the component
in $\mathrm{Sl}(2,\mathbb Z_2)'$ only). We see in

\margin{\uline{N.B.} $g=l_0^{-1}$ is \uline{not}
a solution of this \unsure{the equation|equation}, on the other hand
$g=-l_0^{-1}$ is\ldots}

\scanpage{112}

\unsure{any|any} case from formulas (78) that \unsure{\uline{if}} $\xi$
\deleted{condition}\inserted{is of} the form $\lambda_0\rho(\lambda_0)^{-1}$
\inserted{(if $\xi\rho(\xi)\rho^2(\xi)=1$, of course)},
\deleted{\unsure{then indeed has case A)}} \deleted{\illegible{1}}
\inserted{we will have} $C+D-B=1$ --- i.e. \deleted{\illegible{1}}\inserted{we are}
\inserted{in the} ``general'' case with two parameters.
\deleted{I presume
that if we take
\[
g=\lambda_0(-l_0^{-1})\rho(\lambda_0)^{-1},
\]
we will be in case B), with $c+d-b=-1$
i.e. case 2°) again with two parameters. I
\unsure{dispense with} doing the calculation here.}

\margin{OK, cf. (\unsure{99}) below.}

(Without nevertheless being able to exclude the case
\inserted{where we should at the same time}
$A+B=0$, \deleted{which in fact means that
\[
\xi=\begin{pmatrix}A&B\\C&D\end{pmatrix}
=-\rho^{-1}=\begin{pmatrix}-1&1\\-1&0\end{pmatrix}
\]}
i.e. the case with \unsure{one} parameter
\[
\xi=\begin{pmatrix}-u&u\\C&-u'-C\end{pmatrix}
\]
with $u$ a primitive \unsure{third} root of 1,
$u'=u^{-1}=u^2$ the other primitive root.

\scanpage{113}
\mspage{346}

(VIII) --- \uline{Calculations in} $\mathrm{Sl}(2,\widehat{\mathbb Z})$ : \uline{cocycles under} $\sigma_\infty$.

Let us take again a $g\in\begin{pmatrix}a&b\\c&d\end{pmatrix}\in\mathrm{Sl}(2,\widehat{\mathbb Z})$,
and compute $\sigma_\infty g\sigma_\infty^{-1}$, we find, thanks to
\[
\sigma_\infty=\begin{pmatrix}0&-1\\1&0\end{pmatrix}\qquad (48)
\]
\deleted{\illegible{1}}
the formulas
\[
(90)\quad\left\{
\begin{aligned}
\sigma_\infty g\sigma_\infty^{-1}
&=\begin{pmatrix}d&-c\\-b&a\end{pmatrix}={}^t(g^{-1})
&&\text{the contragredient of $g$ !}\\
&\text{i.e.}\\
\sigma_\infty g^{-1}\sigma_\infty^{\unsure{-1}}
&={}^t g&&\text{thus}\\
[\sigma_\infty,g^{-1}]&=({}^t g)(g)
=\begin{pmatrix}a^2+b^2&ac+bd\\ac+bd&\unsure{c^2+d^2|c^2-d^2}\end{pmatrix}
\end{aligned}\right.
\]
\inserted{$=\sigma_\infty g^{-1}\unsure{\sigma_\infty|\sigma_\infty^{-1}}g$}

\margin{\uline{N.B.} we have
$\sigma_\infty^{-1}=\begin{pmatrix}0&1\\-1&0\end{pmatrix}=-\sigma_\infty$
and \deleted{\illegible{1}} $\sigma_\infty^2\ne1$ in $\mathrm{Sl}(2,\widehat{\mathbb Z})$ !}

\margin{\uline{N.B.} If $g\in\mathrm{Gl}(2,A)$,
$A$ any ring
(for example $A=\mathbb Z_\ell$), we have
$\sigma_\infty g\sigma_\infty^{-1}=\begin{pmatrix}d&-c\\-b&a\end{pmatrix}
=(\det g){}^t(g^{-1})$}

thus, taking the inverse of both sides in
the last formula
\[
(91)\quad \underbrace{g^{-1}\sigma_\infty(g)}_{\inserted{({}^t g\cdot g)^{-1}}}
=\begin{pmatrix}c^2+d^2&-(ac+bd)\\-(ac+bd)&a^2+b^2\end{pmatrix}.
\]

On the other hand, let
\deleted{$h_0=\begin{pmatrix}\unsure{r}&s\\t&u\end{pmatrix}$}
$h_0\in\mathrm{Sl}(2,\widehat{\mathbb Z})$
we have
\[
h_0\sigma_\infty(h_0)=h_0({}^t h_0^{-1})
\]
thus
\[
(*)\qquad h_0\sigma_\infty(h_0)=1\quad\Longleftrightarrow\quad h_0={}^t h_0.
\]

Finally, let
\[
(91)\qquad \xi=\begin{pmatrix}A&B\\C&D\end{pmatrix}
\]
and consider the equation (68)
\[
(**)\qquad \sigma_\infty(\xi)=l_1^\gamma\xi^{-1}l_0^\gamma
\]
which is also equivalent to (*), with $h_0$ given by (\unsure{65|66})
\deleted{$h_0=l_0^{-\unsure{\gamma|\delta}}\xi$ \illegible{4}}, $l_0,l_1$ by (44).

\scanpage{114}

\[
(92)\quad\begin{aligned}
h_0=l_0^{-\gamma}\xi
&=\begin{pmatrix}1&+2\gamma\\0&1\end{pmatrix}
\begin{pmatrix}A&B\\C&D\end{pmatrix}\\
&=\begin{pmatrix}A+2\gamma C&B+2\gamma D\\C&D\end{pmatrix}
\end{aligned}
\]
thus ``the equation in $\xi$'' (**) (written in the
quotient $\mathrm{Sl}(2,\widehat{\mathbb Z})'$ of $\widehat{\mathfrak S}_{0,3}^{+}$, is written $B+2\gamma D=C$
i.e.
\[
(93)\qquad \underbrace{C-B}_{\inserted{1-D}}=2\gamma D
\qquad\text{i.e.}\quad\boxed{pD=1}\ \inserted{\text{\unsure{in case|one case}}}
\]
\deleted{\illegible{2}}

\margin{on \unsure{recalling} that $2\gamma+1=p$}

(We should keep in mind that the equation (**) \uline{in}
$\mathrm{Sl}(2,\widehat{\mathbb Z})'$ means $h_0\sigma_\infty(h_0)=\pm1$ i.e. $h_0=\pm{}^t h_0$,
\deleted{thus here} but the sign $-$ would imply that we have
$D=-D$ i.e. $2D=0$ so $D=0$, which is impossible
if we assume that $\xi=\begin{pmatrix}A&B\\C&D\end{pmatrix}$ comes
from $\widehat\Pi_{0,3}$ so is in $\Pi'(2)$ i.e.
$\equiv\deleted{\unsure{\pm}}1\ (2)$ (so that $D\equiv1\ (2)$).
Thus we must indeed have (93).

Let us now take the parametric
representation \inserted{(70)} of $\xi$ (taking account of
the equation $\xi\rho(\xi)\rho^2(\xi)=1$ \deleted{\unsure{$\xi$}} (64))
\[
(94)\qquad\begin{aligned}
\xi&=\lambda\rho(\lambda)^{-1}
&&\inserted{(\lambda\in\mathrm{Sl}(2,\widehat{\mathbb Z}))}\\
&\text{or}\quad
\xi=\lambda'(\unsure{-l_1^{-1}})\rho(\lambda')^{-1}
=\lambda\rho(\lambda)^{-1}
&&\inserted{\left(\begin{gathered}\lambda=\lambda'\tau_\infty\\\lambda'\in\mathrm{Sl}(2,\widehat{\mathbb Z})\end{gathered}\right)}
\end{aligned}
\]
\deleted{\unsure{$\lambda,\lambda'\in\mathrm{Sl}$}}
with $\lambda\in\mathrm{Gl}(2,\widehat{\mathbb Z})^{\unsure{*|\circ}}$ \inserted{\deleted{\unsure{case I}} $\det\lambda=\pm1$},
(in fact $\lambda\equiv\deleted{\unsure{\pm}}1\ (2)$ in the
case that interests us, but no matter\ldots)

\deleted{A) Case $\xi=\lambda\rho(\lambda)^{-1}=[\lambda,\rho]$ with $\lambda=\begin{pmatrix}a&b\\c&d\end{pmatrix}$,}

\scanpage{115}
\mspage{347}

then formula (78) gives $A,B,C,D$ as
quadratic functions in $a,b,c,d$, \deleted{we will have}
\inserted{(paying Attention \unsure{to the} sign $\delta=-1$ in case 2°)}
\deleted{
\[
(95)\quad\left\{
\begin{aligned}
\unsure{B-C|B+C}&=\unsure{\delta(c^2+d^2+cd-\delta)}=D-1\\
D&=\delta(c^2+d^2+cd)
\end{aligned}\right.
\]
\unsure{$D=$}}
thus (93) is written
\deleted{$1-D=2\gamma D$ \unsure{or} i.e.}
\[
(96)\qquad\boxed{(2\gamma+1)D=1\qquad\text{i.e.}\quad pD=1}
\qquad D=\delta(c^2+d^2+cd),
\]
i.e. the multiplier $p$ is obtained as a function
of $\xi\in\widehat\Pi_{0,3}$ by taking the image
of $\xi$ in $\mathrm{Sl}(2,\widehat{\mathbb Z})'$, in the form $\begin{pmatrix}A&B\\C&D\end{pmatrix}$
(there is a sign ambiguity, removed by the
choice of $\lambda$ such that $\xi=\lambda\rho(\lambda)^{-1}$, the image
of $\lambda$ in $\mathrm{Gl}(2,\widehat{\mathbb Z})'$ \unsure{not|in} determining a
matrix (still denoted $\lambda$) of $\mathrm{Gl}(2,\widehat{\mathbb Z})$ except up to a
\uline{sign}, but $\lambda\rho(\lambda)^{-1}$ no longer depends on the
choice of this \uline{sign}\ldots), by the formula
\[
(97)\qquad p=\frac1D\qquad (D=\unsure{\delta}(c^2+d^2+cd))
\]

Let us now note that on reduction modulo 3,
the quadratic form $c^2+d^2+cd=(c+d)^2-cd$
takes only the values 0 and 1 (clear if $c=0$,
if $d=0$, \deleted{or $c+d=0$ then $f(c,d)=-cd=c^2$}
--- so we
find a perfect square in these \deleted{three} cases, otherwise
$c,d=\pm1$, and $f(c,d)=1+1\pm1\equiv1$ or $0$ \deleted{\illegible{1}} (3). Now here
we must find a unit so the value 0 excluded,
so

\scanpage{116}

\margin{\uline{N.B.} We have $D\equiv p\equiv\delta\ (3)$ (98) knowledge of $D$ (or $p$, or $\delta$) modulo $3$ \deleted{\unsure{already|thus} tells} \unsure{tells} us --- which case we have in the parametric \unsure{representation} of $\xi$. Useless now to treat \unsure{these} cases separately --- since we have $\xi=\lambda\rho(\lambda)^{-1}$, $\lambda=\lambda'\tau_\infty$, $\det\lambda=-1$.}

\[
p\equiv\delta\quad(3)
\]
\inserted{(where $\delta=\det\lambda\in\{\pm1\}$, $\xi=\lambda\rho(\lambda)^{-1}$)}
\deleted{\illegible{4} $\xi=\lambda\rho(\lambda)^{-1}$, $\det\lambda=\delta\in\{\pm1\}$, \illegible{3}}
i.e.
\[
\left\{\begin{aligned}
\text{Case }\xi&=\lambda\rho(\lambda)^{-1}\quad(\lambda\in\mathrm{Sl}(2,\widehat{\mathbb Z}))
\Longleftrightarrow p\equiv1\ (3)\quad\deleted{\unsure{g=\lambda_0\rho(\lambda_0)^{-1}}}\Longleftrightarrow\gamma\equiv0\ (3)\\
\text{Case }\xi&=\lambda'l_1^{-1}\unsure{\rho(\lambda)^{-1}|\rho(\lambda')^{-1}}\quad(\lambda'\in\mathrm{Sl}(2,\widehat{\mathbb Z}))
\Longleftrightarrow p\equiv-1\ (3)\Longleftrightarrow\gamma\equiv-1\ (3).
\end{aligned}\right.
\]

B) Case $\xi=\lambda'l_0^{-1}\rho(\lambda')^{-1}=(\lambda'l_0^{-1})(\rho\lambda'^{-1}\rho^{-1})$.

Using the second formula (77) we find
\[
\xi=\begin{pmatrix}a&b\\c&d\end{pmatrix}
\begin{pmatrix}1&2\\0&1\end{pmatrix}
\begin{pmatrix}a-c&c\\a+b-c-d&c+d\end{pmatrix}
=\begin{pmatrix}A&B\\C&D\end{pmatrix}
\]
with
\[
(99)\quad
\left\{\begin{aligned}
A&=3a^2+b^2+3ab-3ac-2ad-bc-bd
& B&=3ac+2ad+bc+bd\\
C&=-(3c^2+3cd+\deleted{\unsure{2}}d^2)+3ac+ad+2bc+\deleted{\unsure{2}}bd
& D&=3c^2+3cd+d^2
\end{aligned}\right.,
\]
\margin{\uline{N.B.} $C+D-B=\unsure{-1|-\delta}$.}

and formula (93) (which expresses ``the equation in $\xi$'' (**), or the one in $h_0$ (*)) is written, \deleted{\unsure{thus again in the form}} taking account of
\[
B-C=-1+D,
\]
in the form \deleted{\unsure{$2\gamma$}}
\[
2\gamma D=-1-D\qquad\text{i.e.}
\]
\[
(100)\quad \boxed{(2\gamma+1)D=-1\qquad\text{i.e.}\quad pD=-1},
\]
\margin{as before we cannot have $h_0=-{}^t h_0$, for we cannot have $D=0$ if we \unsure{assume} \unsure{that} $\xi\equiv1\ (2)$, \unsure{hence} $D\equiv1\ (2)$.}

so here the multiplier $p$ is expressed as a function of the image of $\xi$ in $\mathrm{Sl}(2,\widehat{\mathbb Z})'$, put in the form $\begin{pmatrix}A&B\\C&D\end{pmatrix}$ where $A,B,C,D\in\widehat{\mathbb Z}$ --- and on removing the sign ambiguity \unsure{as before} --- by the formula
\[
(101)\quad p=-\frac1D\qquad(D=\unsure{3c^2+3cd+d^2|3(c^2+3cd+d^2)}).
\]

\scanpage{117}
\mspage{348}

Here it is clear that the only values modulo $3$ \unsure{represented} by the quadratic form $3c^2+3cd+d^2\equiv d^2\ (3)$ are $0$ and $1$, whence
\[
(102)\quad p\equiv-1\ (3)\quad(\text{i.e. }\gamma\equiv-1\ (3))\quad\text{if }\xi=\lambda l_0^{-1}\rho(\lambda)^{-1}
\]
\[
(\Longleftrightarrow g=\lambda_0l_0^{-1}\rho(\lambda_0)^{-1})
\]

This study thus does show that we are in the case $\xi=\lambda\rho(\lambda)^{-1}$, resp. $\xi=\lambda l_0^{-1}\rho(\lambda)^{-1}$, according as $p=\chi(u)\equiv1$ or $\equiv-1\ (3)$, i.e. according as $\gamma\equiv0$ or $\gamma\equiv-1\ (3)$.

\uline{Remarks.} 1°) \unsure{Placing} ourselves in $\mathrm{Sl}(2,\widehat{\mathbb Z})'$ \deleted{or else} \inserted{(or rather, in its} \deleted{of the $\unsure{\Gamma|F}$} subgroup $\Pi(2)'=\operatorname{Ker}(\mathrm{Sl}(2,\widehat{\mathbb Z})'\unsure{-|\longrightarrow}\mathrm{Sl}(2,\mathbb Z/2\mathbb Z))$ in place of $\widehat{\Pi}_{0,3}$, to \deleted{\unsure{solve the equation}} \inserted{determine the} solutions of ``the equation in $\xi$'' (**) (cf. (68)), when $\xi$ is put in \unsure{parametric} form (94), we see that as a condition on $\lambda=\begin{pmatrix}a&b\\c&d\end{pmatrix}$, this equation is relatively innocuous and very simple --- it simply says that $pD=1$, where $D=c^2+d^2+cd$ in the first case, $D=\deleted{-(3c^2+3cd+d^2)}\inserted{-(c^2+d^2+cd)}$ in the second. If we do not fix $p\in\widehat{\mathbb Z}^{*}$ beforehand, this condition is simply that the element

\scanpage{118}

$D$ of $\widehat{\mathbb Z}$ be a unit, i.e.
\[
(103)\quad D\in\widehat{\mathbb Z}^{*}
\]
--- and we then determine $p$ uniquely in terms of $D$ \inserted{, hence of $c$ and $d$ and the sign $\delta$ by $D=\delta(c^2+d^2+cd)$, \deleted{and of $\delta\in\{\pm1\}$}} ($a$ and $b$ are of no use) --- hence also in terms of $\lambda$\inserted{$\in\widehat{\Pi}_{0,3}\cdot\unsure{(1,\tau)|\{1,\tau\}}$}, itself uniquely determined by $\xi$.
\margin{\uline{N.B.} We \unsure{will have} $p$}

\uline{Thus knowledge of} $\xi\in\widehat{\Pi}_{0,3}$ \uline{suffices} \uline{to determine} \inserted{\uline{the multiplier}} $p=2\gamma+1$\inserted{$\in\widehat{\mathbb Z}^{*}$}, hence also $\gamma\in\widehat{\mathbb Z}$, and $h_0=l_0^{-\gamma}\xi$, hence also $g_0$ (by formula (66), $g_0^{-1}\sigma_\infty(g_0)=h_0$), and consequently $g$ by formula (62), $g=\sigma_\infty(g_0)l_1^\gamma\rho(g_0)^{-1}$, hence $u$ by (50). It is \unsure{interesting} to make sure that \unsure{$D$|it} \deleted{\illegible{1}} does not change when we make the change (69)
\[
\xi\longmapsto\xi'=l_0^{-\alpha}\xi l_1^\alpha,
\]
in matrix terms, we have indeed
\[
(104)\quad l_0^{-\alpha}\xi l_1^\alpha=
\begin{pmatrix}1&2\alpha\\0&1\end{pmatrix}
\begin{pmatrix}A&B\\C&D\end{pmatrix}
\begin{pmatrix}1&0\\2\alpha&1\end{pmatrix}
=\begin{pmatrix}A+2\alpha(B+C)+4\alpha^2D&B+2\alpha D\\C+2\alpha D&D\end{pmatrix}.
\]

2/ But we should take care that in \unsure{full} \uline{rigor} mere knowledge of \uline{the image of $\xi$ in} $\mathrm{Sl}(2,\widehat{\mathbb Z})'$ allows us to determine the multiplier $p$ \uline{only up to sign} --- on the other hand knowledge of the image of $\lambda$ in $\mathrm{Gl}(2,\widehat{\mathbb Z})'$ --- hence of a matrix $\begin{pmatrix}a&b\\c&d\end{pmatrix}\in\unsure{\mathrm{Sl}|\mathrm{Gl}}(2,\widehat{\mathbb Z})$ defined up to \uline{sign} --- allows us to determine

\scanpage{119}
\mspage{349}

$p$ as the inverse of $D$, given \deleted{depending on the case} by $D=\delta(c^2+cd+d^2)$ \deleted{\unsure{or} $D=3(c^2+cd)+d^2$}. \inserted{where $\delta\in\{\pm1\}$ is determined by the condition $D\equiv\delta\ (3)$.}

3/ We know by arithmetic means (since $\Pi_{\mathbb Q}\xrightarrow{\unsure{\chi}}\widehat{\mathbb Z}^{*}$ is surjective) that we must find any $p\in\widehat{\mathbb Z}^{*}$ in the form $1/D$ (or again, in the form $D$), where $D=c^2+cd+d^2$ if $p\equiv1\ (3)$, and $D=-3(c^2+cd)\unsure{-|+}d^2$ if $p\equiv-1\ (3)$, with $c,d\in\widehat{\mathbb Z}$, and even $c\equiv0\ (2)$, $d\equiv1\ (2)$. \inserted{If we want the direct verification,} (We are reduced to verifying it separately for the various $\mathbb Z_\ell$, i.e. to proving the

\uline{Corollary.} \unsure{1°)|°)} If $\ell$ \deleted{\illegible{1}} prime $\ne3$, then $\forall\unsure{\beta|p}\in\mathbb Z_\ell^{*}$ can be put equally well in the forms $c^2+cd+d^2$, or $-(3(c^2+cd)+d^2)$ (with $c,d\in\unsure{\widehat{\mathbb Z}_\ell|\mathbb Z_\ell}$). If $\ell=2$, we can even suppose $c\equiv0$, $d\equiv1\ (2)$. For $p\in\mathbb Z_\ell^{*}$, $\ell=3$, it can be put in the form $c^2+cd+d^2$ (resp. $-3(c^2+cd)\unsure{-|+}d^2$) iff $p\equiv1\ (3)$ (resp. $p\equiv-1\ (3)$), cases which are mutually exclusive.
\margin{we don't care \unsure{about the|about} $-(3(c^2+cd)+d^2)$ !}

4/ \unsure{Returning} to the ``innocuous'' reduction in $\mathrm{Sl}(2,\widehat{\mathbb Z})'$ of ``the equation in $\lambda$'' as meaning simply that $c^2+d^2+cd\in\widehat{\mathbb Z}^{*}$ \deleted{resp. $3(c^2+cd)+d^2\in\widehat{\mathbb Z}^{*}$}, we might hope

\scanpage{120}

to have a method for constructing even in $\widehat{\Pi}_{0,3}$ big batches of solutions of the equation in question. But I have absolutely not managed it --- not even to construct a single one explicitly, apart from \unsure{those} which correspond to the case $u=1$, and $u=\tau_\infty$.

A question: Let $\lambda_0\in\mathrm{Sl}(2,\widehat{\mathbb Z})^{\unsure{\prime|*}}$, $\lambda_0\equiv1\ (2)$, be such that $\xi_0=\lambda_0\rho(\lambda_0)^{-1}$ (say, or $\xi=\lambda(-l_0^{-1})\rho(\lambda_0)^{-1}$) \deleted{be such that} \unsure{satisfies} in $\mathrm{Sl}(2,\widehat{\mathbb Z})'$ an ``equation in $\xi$'' for a suitable $\gamma=\frac{p-1}{2}$ ($p\in\widehat{\mathbb Z}^{*}$) --- i.e. $c^2+d^2+cd\in\widehat{\mathbb Z}^{*}$ (resp. $3(c^2+cd)+d^2\in\widehat{\mathbb Z}^{*}$). Does there then exist a $\unsure{\lambda|\xi}\in\widehat{\Pi}_{0,3}$ above $\lambda_0$, such that $\xi=\lambda\rho(\lambda)^{-1}$ (resp.\ldots) \unsure{satisfies} an equation in $\xi$ as well (\unsure{necessarily} with the same $\gamma$) ?

\scanpage{121}
\mspage{350}

\uline{43} \deleted{Return to the notations} -- Digression on \inserted{the 1-}cocycles of cyclic groups.

\unsure{Let|Let} $\Pi$ be a group, $\rho$ an \unsure{automorphism} of $\Pi$ such that
\[
\rho^n=\mathrm{id} \tag{1}
\]
where $n\in\mathbb N^*$. We thus have an operation of the cyclic group $\mathbb Z/n\mathbb Z$ on $\Pi$, whence a semidirect product
\[
\widetilde\Pi=\Pi.\,\mathbb Z/n\mathbb Z, \tag{\unsure{1|2}}
\]
and an extension structure

(3) \diagram{\begin{tikzcd}[ampersand replacement=\&]
1 \arrow[r] \& \Pi \arrow[r] \& \widetilde\Pi \arrow[r] \& \mathbb Z/n\mathbb Z \arrow[r] \arrow[l,bend left=30,dotted] \& 1.
\end{tikzcd}}

We \deleted{identify}\inserted{denote by} $\dot\rho$ \deleted{\unsure{with}}\inserted{the} generator of $\mathbb Z/n\mathbb Z$, and \unsure{we|to} \inserted{\unsure{we identify $\rho$ with}}\deleted{\illegible{1}} the image of the latter in $\widetilde\Pi$, by the canonical splitting \deleted{\illegible{1}}
\[
\mathbb Z/n\mathbb Z\longrightarrow\widetilde\Pi\qquad 1\bmod n\longmapsto\rho. \tag{4}
\]
We thus have
\[
\rho(g)=\rho g\rho^{-1}\qquad(g\in\Pi). \tag{5}
\]
The splittings of the extension (3) correspond bijectively to the elements $\rho'$ of $\widetilde\Pi$ over $\dot\rho$, such that
\[
\rho'{}^n=1. \tag{6}
\]
To say that $\rho'$ is over $\unsure{\rho|\dot\rho}$ means that we have
\[
\rho'=\rho'_g\overset{\unsure{\mathrm{dfn}}}{=}g\rho\qquad\text{with }g\in\Pi, \tag{7}
\]
and in terms of such a $g\in\Pi$ equation (6) \unsure{Is|is} rewritten as follows
\[
\boxed{g\rho(g)\rho^2(g)\cdots\rho^{n-1}(g)=1} \tag{8}
\]

\scanpage{122}

\deleted{Let us set, for $g\in\Pi$}

\deleted{We will have, for $\alpha\in\Pi$,}
\deleted{$\unsure{\alpha\rho\alpha^{-1}=}\rho'_g=g\rho$}

\deleted{so that \illegible{2}}

\deleted{We will have, for \illegible{1} $g$ as above.}

The group $\Pi$ operates by conjugation on the \unsure{set} of splittings of the extension (3), or what amounts to the same thing, on the \unsure{set} of $\rho'$ above $\unsure{\rho|\dot\rho}$, satisfying (6), \deleted{We} We will moreover have, for $\alpha\in\Pi$
\[
\begin{aligned}
\alpha\rho'_g\alpha^{-1}
&=\deleted{\unsure{(\alpha g\alpha^{-1})(\alpha\rho\alpha^{-1})=\alpha}}\inserted{\alpha g\rho\alpha^{-1}=\alpha g\rho\alpha^{-1}\rho^{-1}\rho}\\
&=(\alpha g\rho(\alpha)^{-1})\rho
=\unsure{\rho'_{\alpha g\rho(\alpha)^{-1}}}
\end{aligned}\tag{9}
\]
i.e. the operation by conjugation of $\Pi$ on the \unsure{set} of splittings of (3), is interpreted as the operation $\Theta=\Theta_\rho$ of $\Pi$ on the \unsure{set} of solutions of (8) \inserted{given, for $\alpha\in\Pi$, by}
\[
\unsure{\Theta_\rho(\alpha)}:g\longmapsto\alpha g\rho(\alpha)^{-1} \tag{10}
\]
(We will say that $\alpha g\rho(\alpha)^{-1}$ is the $\rho$-conjugate of $g$ by the element $\alpha\in\Pi$). The $\Pi$-conjugacy classes \unsure{of|of the} splittings of (3) thus correspond to the $\Pi$-$\rho$-conjugacy classes of solutions of (8).

Among these classes, the \inserted{``trivial''} class \deleted{\unsure{of the extension}}\inserted{\unsure{of the splitting}} canonical for (3), i.e. the class of $g=1$, is given parametrically \deleted{by} as the \unsure{set} of
\[
g=\alpha\rho(\alpha)^{-1}\qquad(\alpha\in\Pi). \tag{11}
\]
Let $g,g'$ be two \unsure{congruent|conjugate} solutions of

\scanpage{123}
\mspage{351}

the equation (8), so that we have
\[
g'=\alpha g\rho(\alpha)^{-1}\qquad\text{i.e. }g'=\Theta(\alpha).g \tag{12}
\]
for a suitable $\alpha\in\Pi$. When do we have uniqueness \unsure{of the|of} $\alpha$ -- more generally, what can we say, for fixed $g,g'$, of the \unsure{set} of $\alpha\in\Pi$ \deleted{which} such that we have (12) (which we can call the $\rho$-\uline{transporter} of $g$ into $g'$, in the group $\Pi$). \deleted{As}\inserted{It is} obviously a right coset \unsure{modulo|module} the subgroup $\Pi(\rho,g)$ of $\Pi$ which \deleted{stabilizes} stabilizes $g$
\[
\begin{aligned}
\Pi(\rho,g)&=\{\beta\in\Pi\mid\Theta_\rho(\beta).g=g\}\\
&=\{\beta\in\Pi\mid\beta g\rho(\beta)^{-1}=g\}\\
&=\{\beta\in\Pi\mid\unsure{\{\mathrm{int}(\beta)\}}(g\rho)=g\rho\}\\
&=\Pi^{\deleted{g}\rho'_g}=\{\beta\in\Pi\mid\mathrm{int}(\underset{\inserted{\rho'_g}}{\uline{g\rho}})\beta=\beta\}\\
&=\text{Centralizer of }\deleted{g}\rho'_g\ \uline{\text{in }\Pi}.
\end{aligned}\tag{13}
\]

In other words, if $\alpha$ satisfies (12), the \unsure{set} of $\alpha'$ which satisfy the same relation is none other than the \unsure{set} $\alpha\Pi(\rho,g)$ :
\[
\begin{aligned}
\alpha'g\rho(\alpha')^{-1}=\alpha g\rho(\alpha)^{-1}
&\Longleftrightarrow\alpha'\in\alpha\Pi(\rho,g)=\alpha\Pi^{\rho'_g}\\
&\text{i.e. }\alpha^{-1}\alpha'\in\Pi^{\rho'_g}
\end{aligned}\tag{14}
\]
Thus we have uniqueness of the $\alpha$ satisfying (12) \deleted{\unsure{iff}} \unsure{iff}
\[
\Pi^{\rho'_g}=1, \tag{15}
\]
a condition which depends only on the conjugacy class of the splitting \deleted{\unsure{under consideration}} of (3) defined by $g$ i.e. by $\rho'_g$.

\scanpage{124}

In fact, not only the group $\Pi$, but the group $\widetilde\Pi$ itself operates by conjugation on the \unsure{set} of splittings of the extension (3), This is in particular \inserted{so} for $\rho$, and we have
\[
\rho(\rho'_g)\overset{\unsure{\mathrm{dfn}}}{=}
\rho\rho'_g\rho^{-1}=\rho(g\rho)\rho^{-1}
=(\rho g\rho^{-1}).\rho=\unsure{g'_{\rho(g)}|\rho'_{\rho(g)}}, \tag{16}
\]
i.e. the operation of $\rho$ on the \unsure{set} of splittings in question, is identified with the operation \deleted{of} induced \deleted{by}\unsure{by|$\rho$} the permutation $\rho$ of $\Pi$ on the \unsure{set} of solutions of (8) $(g\in\Pi)$. Let us note moreover that $g$ and $\rho(g)$ \deleted{\illegible{1}} are $\Pi$-conjugate, for
\[
\rho(g)=\deleted{g}\inserted{\alpha}g\rho(\alpha)^{-1}\qquad\text{with }\alpha=g^{-1}. \tag{17}
\]

Let us note that \deleted{\illegible{4}}\inserted{the operation of $\widetilde\Pi$}\deleted{\unsure{$\rho$}}, on the \unsure{set} of splittings of (3), or again on the \unsure{set} of solutions of (8), is \unsure{always|thus} \unsure{given} by
\[
\Theta_\alpha:g\longmapsto\alpha g\rho(\alpha)^{-1} \tag{10}
\]
where now we assume $\alpha\in\widetilde\Pi$ (and not necessarily $\alpha\in\Pi$). The $\widetilde\Pi$-conjugacy classes are \unsure{reduced|identical} to the $\Pi$-conjugacy classes, \inserted{(as it follows} \unsure{follows} from (17)) \deleted{\unsure{has just seen it.}} In particular, if $\alpha\in\widetilde\Pi$, the element
\[
\alpha\rho(\alpha)^{-1}\qquad(\alpha\in\widetilde\Pi) \tag{18}
\]
of $\widetilde\Pi$ is in $\Pi$, and it \deleted{\illegible{2}} satisfies (8), and

\scanpage{125}
\mspage{352}

is in the principal class, i.e. we will have
\deleted{\unsure{(19)}}
\[
\alpha\rho(\alpha)^{-1}=\alpha_0\rho(\alpha_0)^{-1}\qquad\text{for }\alpha_0\in\Pi\text{ suitable}
\]
In fact, we can write $\alpha=\alpha_0\rho^i$\inserted{, with $\alpha_0\in\Pi$, $i\in\mathbb Z$}, whence $\rho(\alpha)^{-1}=\rho^{-i}\rho(\alpha_0)^{-1}$ whence (19). More generally, we will have, \deleted{\unsure{if $\alpha,\alpha_0$}}
\[
\alpha=\alpha_0\rho^i\qquad(\deleted{\unsure{\alpha\in\widetilde\Pi}}\alpha_0\in\Pi,\ i\in\mathbb Z)
\]
the relation
\deleted{\unsure{(19)}}
\[
\begin{aligned}
\alpha g\rho(\alpha^{-1})
&=\alpha_0\underset{\deleted{\unsure{\rho^i(g)}}}{\underbrace{\rho^i g\rho^{-i}}}\unsure{\rho(\alpha_0^{-1})|\rho(\alpha_0)^{-1}}\\
&=\deleted{\unsure{\alpha_0\rho^i g\rho(\alpha_0\rho^i)^{-1}}}\inserted{\alpha_0\rho^i(g)\rho(\alpha_0)^{-1}}
\end{aligned}
\]
\unsure{with|for}
\[
\rho^i(g)=\beta g\rho(\beta)^{-1}\qquad
\beta=(g\rho(g)\cdots\rho^{i-1}(g))^{-1}
\]
thus
\[
\left\{\begin{aligned}
\alpha g\rho(\alpha^{-1})&=(\alpha_0\beta)g\rho(\alpha_0\beta)^{-1}\\
\beta&=(g\rho(g)\cdots\rho^{i-1}(g))^{-1}
\end{aligned}\right. \tag{19}
\]
(for $\alpha=\alpha_0\rho^i$, $i\in\unsure{\mathbb N|\mathbb N^*}$).

Of course, the representation (18) of a solution of equation (8), \deleted{\unsure{is not}} when we no longer assume $\alpha\in\Pi$ but only $\alpha\in\widetilde\Pi$, is no longer one-to-one, since for $\alpha=\rho$ for example we find \deleted{\unsure{the element 1}} $g=\rho\rho^{-1}=1$.

\deleted{Let us suppose \unsure{in} addition}\inserted{Let us suppose} more generally that we have a group
\deleted{\unsure{(20)}}
\[
\unsure{\mathcal E|\varepsilon}\supset\widetilde\Pi
\]
(as in the \deleted{\unsure{group}} case where $\Pi=\Pi_{0,3}$, $\widetilde\Pi$ inverse image of $\{1,\dot\rho,\dot\rho^2\}\subset\mathfrak S_3$ in $\mathfrak S_{0,3}$ extension of $\mathcal T_{0,3}$ by $\Pi_{0,3}$), in which $\Pi$ is normal, so that \deleted{$\unsure{\mathcal E|\varepsilon}$ operates by conjugation}\inserted{we have an extension}

\scanpage{126}

\deleted{on the \unsure{set} of $\rho'$ in $\Pi$ \unsure{such|such} that}

\[
(20)\qquad 1\longrightarrow\Pi\longrightarrow\unsure{\mathcal E|\varepsilon}\longrightarrow G\longrightarrow1
\]

\margin{N.B. the extension (3) is the inverse image of (20) under the inclusion $\mathbb Z/n\mathbb Z\hookrightarrow G$}

\unsure{with}
\[
(21)\qquad \mathbb Z/n\mathbb Z\hookrightarrow G,\qquad
1\bmod n\longmapsto\dot\rho\in G
\]
\inserted{(element of order $n$)}
$\widetilde\Pi$ being identified with the inverse image of $\mathbb Z/n\mathbb Z$ in $\unsure{\mathcal E|\varepsilon}$.

We cannot say in general, if $\rho'$ is an \unsure{element} of $\unsure{\mathcal E|\varepsilon}$ above $\dot\rho$ (thus $\rho'\in\widetilde{\unsure{\mathcal E|\varepsilon}}$).
(even if it satisfies $\rho'{}^n=1$), that for $\forall$ $\alpha\in\unsure{\mathcal E|\varepsilon}$, $\operatorname{int}(\alpha)(\rho')$ is still above $\dot\rho$:
for this to be so, it \unsure{is necessary and sufficient} that $\dot\alpha\in G$ centralize $\dot\rho$. So let $\unsure{\mathcal E|\varepsilon}_\rho$
\[
(22)\qquad
\unsure{\mathcal E|\varepsilon}_\rho=\{\alpha\in\unsure{\mathcal E|\varepsilon}\mid[\dot\alpha,\dot\rho]=1\}
\]
it is a subgroup of $\unsure{\mathcal E|\varepsilon}$ containing $\widetilde\Pi$, which
\deleted{\unsure{for $\forall$}}\inserted{\unsure{itself}} operates on the set of \deleted{\unsure{$\rho'$}} $\rho'$ above $\dot\rho$,
\inserted{in particular on \deleted{\illegible{1}} the set of \illegible{1} \unsure{which}}
\unsure{satisfying} (8), \deleted{and it} by the same operation
\[
(23)\qquad
\alpha\rho'_g\alpha^{-1}=\rho'_{\alpha g\rho(\alpha)^{-1}}
\]
so that we can also consider this \unsure{operation} as given by
\[
(24)\qquad
\left\{
\begin{array}{l}
\Theta_{\rho,\unsure{\mathcal E|\varepsilon}_\rho}(\alpha)=\Theta(\alpha):
g\longmapsto\alpha g\rho(\alpha)^{-1}\\
(\alpha\in\unsure{\mathcal E|\varepsilon}_\rho).
\end{array}
\right.
\]

This time, an operation $\Theta(\alpha)$ may permute the $\Pi$-conjugacy classes of the \deleted{\unsure{various}} solutions, of (8) in $\Pi$.
\unsure{We|It} \deleted{\illegible{1}} will have on $\unsure{\mathcal E|\varepsilon}_\rho$

\scanpage{127}
\mspage{353}

an extension structure
\[
(25)\qquad
\begin{array}{ccccccccc}
1&\longrightarrow&\Pi&\longrightarrow&\unsure{\mathcal E|\varepsilon}_\rho&\longrightarrow&G^\rho&\longrightarrow&1\\
&&&&&&\Vert\\
&&&&&&\multicolumn{3}{l}{\operatorname{Centr}_G(\rho)=\{g\in G\mid\rho g\unsure{g|\rho}^{-1}=g\}}\\
&&&&&&\multicolumn{3}{l}{\phantom{\operatorname{Centr}_G(\rho)}=\{g\in G\mid[\rho,g]=1\}}
\end{array}
\]
and $\unsure{\mathcal E|\varepsilon}_\rho$ operates on the set of $\Pi$-conjugacy classes of splittings of (3), (or the $\Pi$-$\rho$-conjugacy classes of solutions of (8)) via an operation of $G^\rho$ on it.

Let $\deleted{\unsure{g}}\alpha\in\unsure{\mathcal E|\varepsilon}$, let us consider the element
\[
\deleted{\unsure{g}}\alpha\rho(\deleted{\unsure{g}}\alpha)^{-1}
=\alpha(\rho\alpha\rho^{-1})^{-1}
=\alpha\rho\alpha^{-1}\rho^{-1}=[\alpha,\rho]
\]
It is in $\Pi$ iff we have \deleted{\unsure{that}}
\[
[\dot\alpha,\dot\rho]=1\qquad\text{i.e.}\quad
\dot\alpha\in G^\rho\qquad\text{i.e.}\quad\alpha\in\unsure{\mathcal E|\varepsilon}_\rho.
\]

For a solution $g$ of equation (8), it may happen that we cannot put \unsure{it|it} in the \inserted{standard} form $\alpha\rho(\alpha)^{-1}$ with $\alpha\in\Pi$, but that we succeed with $\alpha\in\unsure{\mathcal E|\varepsilon}_\rho$.
For such an $\alpha$ to exist, it is necessary and sufficient that the $\Pi\rho$-conjugacy class \inserted{$[g]_\rho$} of $g$ (among the solutions of (8)) be in the orbit of the trivial conjugacy class \inserted{$[1]_\rho$}, for the action under consideration of $G^\rho$ on this set of conjugacy classes; \deleted{and}\inserted{then} the $\alpha\in\unsure{\mathcal E|\varepsilon}_\rho$ which \unsure{satisfy}
\[
(26)\qquad g=\alpha\rho(\alpha)^{-1}
\]

\scanpage{128}

\deleted{are such that}
\deleted{satisfy}
\deleted{\unsure{are some} \illegible{1}}
\inserted{\unsure{satisfy} the relation}
\[
(27)\qquad\dot\alpha[1]_\rho=[g]_\rho,
\]
and for $\forall\ \dot\alpha\in G^\rho$ satisfying this relation (27), there exists at least one representative $\alpha\in\unsure{\mathcal E|\varepsilon}_\rho$ satisfying (6) (which is \unsure{pure} tautology).

Under what condition do we have \uline{uniqueness} \unsure{of the|of the} $\alpha\in\unsure{\mathcal E|\varepsilon}_\rho$ satisfying (26)?
The set of these $\alpha$ is a right coset under the \unsure{subgroup} of $\unsure{\mathcal E|\varepsilon}_\rho$ \unsure{formed} of the $\alpha$ such that \deleted{\illegible{1}}
\[
\alpha\rho(\alpha)^{-1}=1\qquad\text{i.e.}\quad
\alpha=\rho(\alpha)\qquad\text{i.e.}\quad[\alpha,\rho]=1,
\]
this group is none other than \deleted{$\unsure{\mathcal E|\varepsilon}$}
\[
(28)\qquad\unsure{\mathcal E|\varepsilon}^\rho=\operatorname{Centr}_{\unsure{\mathcal E|\varepsilon}}(\rho)\subset\unsure{\mathcal E|\varepsilon}_\rho.
\]
\unsure{which} group contains $\rho$ \unsure{hence} the \inserted{\unsure{sub-}}group $\simeq\mathbb Z/n\mathbb Z$ \inserted{of $\unsure{\mathcal E|\varepsilon}$} generated by $\rho$ -- so it is never trivial -- i.e. we never have uniqueness (except in the trivial case where \deleted{\unsure{$\rho$}} $n=1$ i.e. $\rho=1$, \deleted{\illegible{1}} but then $\unsure{\mathcal E|\varepsilon}^\rho=\unsure{\mathcal E|\varepsilon}$ and we would need $\unsure{\mathcal E|\varepsilon}=1$ \dots !)

But if we assume $\Pi^\rho=\{1\}$ i.e.
\[
(29)\qquad\unsure{\mathcal E|\varepsilon}^\rho\cap\Pi=\{1\}
\]
\deleted{\unsure{then}}
\[
\unsure{\mathcal E|\varepsilon}^\rho_{\deleted{\unsure{\rho}}}
\underset{\uline{\text{injective}}}{\overset{\mathrm{can}}{\hookrightarrow}}
G^\rho\subset G,
\]
then the $\alpha\in\unsure{\mathcal E|\varepsilon}^\rho$ solutions of (26) are in one-to-one correspon-

\scanpage{129}
\mspage{354}

dence with their images \inserted{$\dot\alpha$} in $G$, i.e. with the transporter \inserted{in $G$} from $[1]_\rho$ \unsure{into|to} $[g]_\rho$.
If we choose an element of this transporter, we have chosen by \unsure{that very fact} an $\alpha$, solution of (26).

\[
\boxed{\begin{gathered}
\text{In the case of }\quad
\Pi=\Pi_{\unsure{0,3}}\inserted{=\Pi_1(U_{0,3})},\qquad
\unsure{\mathcal E|\varepsilon}=\mathfrak S_{0,3}
\inserted{=\Pi_1(U_{0,3},\mathcal T_{0,3})},
\quad\deleted{\unsure{G=\mathcal T_0}}\\
G=\mathcal T_{0,3}\simeq\mathfrak S_3\times\mathbb Z/2\mathbb Z
\end{gathered}}
\]
with $\rho,\dot\rho$ as usual, the set of conjugacy classes of the \deleted{\unsure{solutions}}\inserted{splittings} of (3) i.e. of the solutions of (8), \unsure{is identified} with $U_{0,3}^{\dot\rho}$, which has two elements, \unsure{the} two ``poles'' $-\overline j,-j$ -- one of which, chosen as base point, corresponds to the trivial class.

\deleted{\unsure{Here the} \illegible{1} $G^\rho\simeq\unsure{\mathbb Z/2\mathbb Z}\{1,\tau\}\simeq\mathbb Z$ is the group}
\inserted{Here}
\[
G^\rho=\unsure{\mathcal T|\mathcal G}_{0,3}^{\rho}
\simeq\mathbb Z/3\mathbb Z\times\mathbb Z/2\mathbb Z
\]
is the group generated by $\rho$ and $\tau$, and $\rho$ operates \unsure{trivially|trivially} on the set of the two classes, while $\tau$ permutes them (i.e. $\mathbb Z/\unsure{2|n}\mathbb Z$ \inserted{generated by $\tau$} operates on it in a simply transitive way).
Thus if we \unsure{set}
\[
\begin{aligned}
(30)\qquad\unsure{\mathcal E|\varepsilon}'
&=\inserted{\unsure{\widetilde\Pi|\Pi}_{0,3}=}\Pi\cdot\{1,\tau_\infty\}\\
&=\text{inverse image of }\{1,\tau_\infty\}\subset G\\
&\phantom{={}}\text{in }\unsure{\mathcal E|\varepsilon}=\mathfrak S_{0,3}
\end{aligned}
\]
\margin{(\deleted{\illegible{1}} unoriented \unsure{weighted} triangulated cartographic \unsure{group} \dots}
we see that
\[
(31)\qquad\unsure{\mathcal E|\varepsilon}'{}^\rho=\{1\},
\]
\deleted{(immediate consequence of the relation $\Pi^\rho=(1)$)}
and the set of solutions in \uline{$\Pi_{0,3}$} of the equation
\[
(32)\qquad g\rho(g)\rho^2(g)=1
\]

\scanpage{130}

is in 1-1 correspondence with the set of elements of $\widetilde\Pi_{0,3}=\unsure{\mathcal E|\varepsilon}'$ (the unoriented weighted \inserted{triangulated} cartographic group, having as generators the elements $\tau_0,\tau_1,\tau_\infty$ related by the relations $\tau_0^2=\tau_1^2=\tau_\infty^2=1$ and no more, with $l_0=\tau_\infty\tau_1$, $l_1=\tau_0\tau_\infty$, $l_\infty=\tau_1\tau_0$ \dots)
It should still be likewise when we place ourselves in the groups $\Pi_{0,3}^{\wedge}$ and $\widetilde\Pi_{0,3}^{\wedge}=\widetilde\Pi_{0,3}\cdot\{1,\tau_\infty\}$, if we admit (something to be checked later)
\[
(33)\qquad(\widetilde\Pi_{0,3}^{\wedge})^\rho=\{1\}
\]
\inserted{and that the elements of finite order of $\widehat{\mathfrak S}_{0,3}$ are conjugates of elements of finite order of $\mathfrak S_{0,3}$ \dots}

Let us note that here we have
\[
(34)\qquad\tau_\infty\rho(\tau_\infty)^{-1}
=\tau_\infty\rho(\tau_\infty)=l_1^{-1},
\qquad\deleted{\unsure{\text{in}}}
\]
in virtue of the formulas (23) of \S\ 42
\[
\rho_\infty\overset{\unsure{\mathrm{dfn}|\text{definition}}}{=}
\tau_\infty\rho\tau_\infty^{-1}
=l_1^{-1}\rho=\rho l_0^{-1}.
\]

The situation is even simpler for the \unsure{element|element} $\sigma_\infty$ of $\mathfrak S_{0,3}$, operating on $\Pi$ by an automorphism of order $n=2$, for here every solution of the \unsure{corresponding} equation
\[
(35)\qquad\deleted{\unsure{(36)}}\quad
\deleted{\unsure{g}}h\,\sigma_\infty(\deleted{\unsure{g}}h)=1
\qquad(h\in\Pi)
\]
can be put in the standard form
\[
(36)\qquad h=\beta\sigma_\infty(\beta)^{-1}
\qquad(\beta\in\Pi)
\]
and this in a \uline{unique way}. This \unsure{follows} indeed from the general theory \dots, taking into account that
\[
(37)\qquad U_{0,3}^{\dot\sigma_\infty}=\{\unsure{\tfrac12}\}
\]
